\documentclass[10pt,a4paper]{amsart}
\usepackage[margin=19mm]{geometry}
\usepackage{amsmath,amssymb,mathtools}
\usepackage[scr=rsfs,cal=euler]{mathalfa}
\usepackage{xfrac}
\usepackage[unicode,hidelinks,bookmarks=false]{hyperref}
\newcommand{\ie}{i.e.\ }
\newcommand{\cf}{cf.\ }
\newcommand{\resp}{respectively\ }
\newcommand{\Subsection}{\subsection}
\providecommand{\nobreakdash}{\nobreak\hbox{-}}
\setlength{\emergencystretch}{2em}
\multlinegap0pt
\usepackage{bm}
\usepackage{bbm}
\usepackage{dsfont}
\usepackage{xspace}
\usepackage{cancel}
\usepackage{mathtools}
\usepackage{amsthm}
\makeatletter
\@ifundefined{theorem}{\newtheorem{theorem}{Theorem}}{}
\@ifundefined{lemma}{\newtheorem{lemma}[theorem]{Lemma}}{}
\makeatother
\newcommand{\D}{\mathbb{D}}
\newcommand{\R}{\mathbb{R}}
\newcommand{\T}{\mathbb{T}}
\newcommand{\bx}{\mathbf{x}}
\newcommand{\bz}{\mathbf{z}}
\newcommand{\cB}{\mathcal{B}}
\newcommand{\dist}{\mathrm{dist}}
\newcommand{\trans}{\mathrm{trans}}
\title[Oscillatory integral operators and variable Schr\"odinger pr]{Oscillatory integral operators and variable Schr\"odinger propagators: beyond the universal estimates}
\author{Mingfeng Chen, Shengwen Gan, Shaoming Guo, Jonathan Hickman, Marina Iliopoulou, James Wright}
\date{Source: arXiv:2407.06980v2 (CC BY 4.0)}
\begin{document}
\maketitle

\noindent\emph{Source and licence.} Oscillatory integral operators and variable Schr\"odinger propagators: beyond the universal estimates. Version arXiv:2407.06980v2 (CC BY 4.0), distributed under the Creative Commons Attribution 4.0 International licence, as recorded in the arXiv licence metadata for this exact version and shown on the abstract page https://arxiv.org/abs/2407.06980v2. Full rights notice: licenses/arxiv-2407.06980-CC-BY-4.0.txt.\par\smallskip

\noindent\textit{Editorial scope.} Counted unit: the Lemma whose source label is \texttt{lem: b/n dec pointwise} in the article (source0.tex lines 2746--2752, source label \texttt{lem: b/n dec pointwise}), delivered together with the article's own complete original proof of it (source0.tex lines 2783--2852, one \texttt{proof} environment of 70 lines). No mathematics is written, repaired or re-derived: statement and proof are the article's own text line for line. The proof is detached from the statement in the source: the article states the result at lines 2746--2752 and prints its proof at lines 2783--2852, and the proof environment's own header names the statement, so the association is the article's own. Required context retained uncounted: eq: V aligned caps (lines 2732--2734); lem: transversality (lines 2757--2765); eq: broad points (lines 2758--2760). Every definition, display and label of those lines is retained unchanged. Omitted from this package and not redistributed: 1--2731, 2735--2745, 2753--2756, 2761--2782, 2853--4458. Nothing inside a retained excerpt is omitted or altered: every retained body is delivered exactly as the article's lines are written, so no omission receipt of any kind accompanies this package. Citations: the retained excerpts carry no citation command at all, so no bibliography entry of the article is transcribed and no reference is redistributed. Reference adaptation: none was needed and none was made; every reference inside the delivered unit resolves inside it, so no unresolved reference or citation is produced. Printed slips kept literal: no printed word, symbol, hypothesis or number of the counted statement or of its proof was altered. Layout and class: the article's own class is \texttt{amsart} with packages including inputenc, amssymb, amsmath, amsthm, color, enumerate, mathabx, mathrsfs, tikz, bbm, url, pdfpages, soul, esint; the wrapper is the lane's pinned \texttt{amsart} editorial wrapper at 10pt on A4, which restates the theorem-like environments the retained text needs and adds the article's own macro definitions that the retained text needs (\texttt{\textbackslash D}, \texttt{\textbackslash R}, \texttt{\textbackslash T}, \texttt{\textbackslash bx}, \texttt{\textbackslash bz}, \texttt{\textbackslash cB}, \texttt{\textbackslash dist}, \texttt{\textbackslash trans}), with extra wrapper packages (bm, bbm, dsfont, xspace, cancel, mathtools). The delivered numbering is the unit's own. Assets: no retained span contains a figure environment, a graphics-inclusion command, a PDF-page inclusion, a table or a TikZ picture; no image file is copied into this package, the package declares no asset dependency and no image file is redistributed with it. Licence: version arXiv:2407.06980v2 of this article is distributed under the Creative Commons Attribution 4.0 International licence, as recorded in the arXiv licence metadata for this exact version and shown on the abstract page https://arxiv.org/abs/2407.06980v2. A full rights notice is delivered at licenses/arxiv-2407.06980-CC-BY-4.0.txt. Characters: every retained excerpt is pure ASCII, so the delivered engine, XeLaTeX with its default Latin Modern text font, reports no missing character.
\begin{equation}\label{eq: V aligned caps}
    \Theta(K^{-1};V; \bx) := \big\{ \theta \in \Theta(K^{-1}) : \dist(G(\bx ;y_{\theta}), V) \leq C_{\phi} K^{-1} \big\},
\end{equation}

\begin{lemma}\label{lem: transversality} Let $0 \leq d \leq n-1$ and suppose $\Omega \subseteq \R^n$ is a non-empty subset of the unit sphere $S^{n-1}$ which satisfies the following condition:
\begin{equation}\label{eq: broad points}
    \textrm{for all $W \in \mathrm{Gr}(d,\R^n)$, there exists some $\omega \in \Omega$ such that $\dist(\omega, W) \geq K^{-1}$.}
\end{equation}
Then there exist $\omega_1, \dots, \omega_{d+1} \in \Omega$ such that 
\begin{equation*}
    |\omega_1 \wedge \cdots \wedge \omega_{d+1}| \geq K^{-d}.
\end{equation*}
\end{lemma}

\begin{lemma}[Broad-narrow decomposition]\label{lem: b/n dec pointwise} Let $2 \leq k \leq n$, $\delta > 0$, $1 \leq K \ll \delta^{-1}$ and $\phi \colon \D^n \to \R$ a non-degenerate phase. Given $\T$ a family of $\delta$-tubes associated to $\phi$ and $\bx \in \R^n$, we have
\begin{align}\label{eq: b/n dec pointwise}
\mu_{\T}(\bx) \lesssim \max_{V \in \mathrm{Gr}(k-1,\R^n)} &\sum_{\theta \in \Theta(K^{-1}; V; \bx)}\mu_{\T_{\mathrm{K}}[\theta]}(\bx) \\
\nonumber
&\qquad + K^{n-1} \max_{(\theta_1, \dots, \theta_k) \in \Theta^k_{\trans}(K^{-1}; \bx)}  \prod_{\ell = 1}^k \mu_{\T_{\mathrm{K}}[\theta_{\ell}]}(\bx)^{1/k}.
\end{align}
\end{lemma}

\begin{proof}[Proof (of Lemma~\ref{lem: b/n dec pointwise})] The first step is to identify some $V_{\bx} \in \mathrm{Gr}(k-1,\R^n)$ which `captures' as many of the large values of $\mu_{\T_{\mathrm{K}}[\theta]}(\bx)$ as possible. More precisely, define the \textit{broad part} of the multiplicity function
\begin{equation}\label{eq: gen b/n 1} 
    \mu_{\T}^{\mathrm{Br}}(\bx) := \min_{V \in \mathrm{Gr}(k-1,\R^n)} \max_{\theta \notin \Theta(K^{-1}; V; \bx)} \mu_{\T_{\mathrm{K}}[\theta]}(\bx)
\end{equation}
and suppose $V_{\bx} \in \mathrm{Gr}(k-1,\R^n)$ realises the minimum in \eqref{eq: gen b/n 1}. For slightly technical reasons, we also define 
\begin{equation}\label{eq: gen b/n 2} 
    \widetilde{\Theta}(K^{-1}; V; \bx) := \big\{ \theta \in \Theta(K^{-1}; V; \bx) : \mu_{\T_{\mathrm{K}}[\theta]}(\bx) \geq \mu_{\T}^{\mathrm{Br}}(\bx) \big\}
\end{equation}
and further choose $V_{\bx}$ so that, of all spaces realising the minimum in \eqref{eq: gen b/n 1}, the space $V_{\bx}$ also maximises $\#\widetilde{\Theta}(K^{-1}; V; \bx)$. 

From the definition, the normals $G(\bx;y_{\theta})$ associated to the caps $\theta \in \Theta(K^{-1}; V_{\bx}; \bx)$ are aligned around the $(k-1)$-dimensional subspace $V_{\bx}$. However, they do not align around any lower dimensional subspace. \medskip

\noindent \textbf{Claim.} There does not exist an affine subspace $W \subseteq \R^n$ of dimension $k-2$ such that $\widetilde{\Theta}(K^{-1}; V_{\bx}; \bx) \subseteq \Theta(K^{-1}; W; \bx)$. \medskip

Assuming the claim, it is a simple matter to conclude the proof of Lemma~\ref{lem: b/n dec pointwise}. We define the collections of \textit{narrow} and \textit{broad} caps (for $\mu_{\T}$ at $\bx$) by
\begin{equation*}
    \mathcal{N}_{\bx} := \Theta(K^{-1}; V_{\bx}; \bx) \qquad \textrm{and} \qquad  \cB_{\bx} := \Theta(K^{-1}) \setminus \Theta(K^{-1}; V_{\bx}; \bx),
\end{equation*}
respectively. By the defining properties of $V_{\bx}$, we have
\begin{align*}
    \mu_{\T}(\bx) &\leq  \sum_{\theta \in \mathcal{N}_{\bx}} \mu_{\T_{\mathrm{K}}[\theta]}(\bx)  + \sum_{\theta \in \cB_x} \mu_{\T_{\mathrm{K}}[\theta]}(\bx)  \\
    &\lesssim \max_{V \in \mathrm{Gr}(k-1,\R^n)} \Big[ \sum_{\theta \in \Theta(K^{-1}; V; \bx)} \mu_{\T_{\mathrm{K}}[\theta]}(\bx) \Big]  + 
 K^{n-1} \mu_{\T}^{\mathrm{Br}}(\bx).   
\end{align*}

In view of the claim, there exists a $(k-1)$-tuple of caps 
\begin{equation*}
(\theta_{\bx,1}, \dots, \theta_{\bx,k-1}) \in \widetilde{\Theta}(K^{-1}; V_{\bx}; \bx)^{k-1}
\end{equation*}
which is transverse at $\bx$ in the sense that
\begin{equation*}
        \inf  \Big\{ \Big|\bigwedge_{\ell = 1}^{k-1} G(\bz;y_{\ell}) \Big| : y_{\ell} \in \theta_{x,\ell}, \, 1 \leq \ell \leq k -1, \, |\bz-\bx|_{\infty} < 2K^{-1} \Big\}\geq K^{-(k-2)}. 
\end{equation*}
To see this, we simply apply Lemma~\ref{lem: transversality} with
\begin{equation*}
    \Omega := \big\{G(\bx; y_{\theta}) : \theta \in \widetilde{\Theta}(K^{-1}; V_{\bx}; \bx)^{k-1}\big\},
\end{equation*}
using the claim to verify the hypothesis \eqref{eq: broad points} with $d=k-2$. The constant $C_{\phi}$ in the definition of \eqref{eq: V aligned caps} can be used to pass from bounds on the the size of wedge products of the $G(\bx; y_{\theta})$ to more general $G(\bz;y_{\theta})$. 

By definition, there exists a cap $\theta_{\bx, k} \notin \Theta(K^{-1}; V_{\bx}; \bx)$ such that 
\begin{equation*}
    \mu_{\T_{\mathrm{K}}[\theta_{\bx,k}]}(\bx) = \mu_{\T}^{\mathrm{Br}}(\bx).
\end{equation*}
It follows that the $k$-tuple $(\theta_{\bx,1}, \dots, \theta_{\bx, k})$ is transversal at $\bx$ and satisfies
\begin{equation*} 
    \mu_{\T}^{\mathrm{Br}}(\bx) \leq \prod_{\ell=1}^k \mu_{\T_{\mathrm{K}}[\theta_{\bx, \ell}]}(\bx)^{1/k} \leq \max_{(\theta_1, \dots, \theta_k) \in \Theta^k_{\trans}(K^{-1}; \bx)}  \prod_{\ell = 1}^k \mu_{\T_{\mathrm{K}}[\theta_{\ell}]}(\bx)^{1/k}.
\end{equation*}
\smallskip

It remains to prove the claim. We argue by contradiction, assuming such a subspace $W$ exists. Let $\theta^* := \theta_{\bx, k}$ be as above, so that $\theta^* \notin \Theta(K^{-1}; V_{\bx}; \bx)$ realises the maximum in the definition \eqref{eq: gen b/n 1} for $V = V_{\bx}$. Define $V_{\bx}^*$ to be a linear subspace of dimension $k-1$ which contains $W$ and $G(\bx; y^*)$, where $y^*$ is the centre of $\theta^*$.

First note that $V_{\bx}^*$ realises the minimum in \eqref{eq: gen b/n 1}. Indeed, by the definition \eqref{eq: gen b/n 2}, if $\theta \in \Theta(K^{-1}; V_{\bx}; \bx) \setminus  \widetilde{\Theta}(K^{-1}; V_{\bx}; \bx)$, then $\mu_{\T_{\mathrm{K}}[\theta]}(\bx) < \mu_{\T}^{\mathrm{Br}}(\bx)$. Consequently,
\begin{equation*}
    \mu_{\T}^{\mathrm{Br}}(\bx) =  \max_{\theta \notin \Theta(K^{-1}; V_{\bx}; \bx)} \mu_{\T_{\mathrm{K}}[\theta]}(\bx) = \max_{\theta \notin \widetilde{\Theta}(K^{-1}; V_{\bx}; \bx)} \mu_{\T_{\mathrm{K}}[\theta]}(\bx).
\end{equation*}
On the other hand, from the hypothesis on $W$ and the definition of $V_{\bx}^*$ we have $\widetilde{\Theta}(K^{-1}; V_{\bx}; \bx) \subseteq \Theta(K^{-1}; V_{\bx}^*; \bx)$ and so
\begin{equation*}
    \mu_{\T}^{\mathrm{Br}}(\bx) \geq \max_{\theta \notin \Theta(K^{-1}; V_{\bx}^*; \bx)} \mu_{\T_{\mathrm{K}}[\theta]}(\bx) \geq \mu_{\T_{\mathrm{K}}[\theta^*]}(\bx) = \mu_{\T}^{\mathrm{Br}}(\bx). 
\end{equation*}

In light of the maximality of $V_{\bx}$, it follows that $\# \widetilde{\Theta}(K^{-1}; V_{\bx}^*, \bx) \leq \# \widetilde{\Theta}(K^{-1}; V_{\bx},\bx)$. However, it is clear from the definitions that 
\begin{equation*}
    \widetilde{\Theta}(K^{-1}; V_{\bx}; \bx) \subseteq \widetilde{\Theta}(K^{-1}; V_{\bx}^*, \bx), 
\end{equation*}
whilst
\begin{equation*}
    \theta^* \in \widetilde{\Theta}(K^{-1}; V_{\bx}^*, \bx) \quad \textrm{and} \quad \theta^* \notin \widetilde{\Theta}(K^{-1}; V_{\bx}; \bx).
\end{equation*}
 Thus, $\# \widetilde{\Theta}(K^{-1}; V_{\bx}; \bx) < \# \widetilde{\Theta}(K^{-1}; V_{\bx}^*, \bx)$, which is a contradiction. 
\end{proof}

\end{document}
